\section{Finite-cover mixing and the hypotheses of the action principle}
\label{sec:v37-cover-inputs}
The arithmetic argument uses an ergodic assertion on actual finite
covers, in addition to local collision geometry. We give its precise
source and keep its constants distinct from the parameter-uniform
constants used on the base table.

\begin{proposition}[Mixing of the physical finite covers]
\label{prop:v37-finite-cover-mixing}
Fix a positive integer $p$. On the torus $\R^2/(p\Lambda)$, the
collision map of each circular or oriented-ellipse table considered
in this article is mixing for normalized collision measure. The same
holds for the support-function tables in
Proposition~\ref{prop:v37-support-family}. For each fixed $p$, the
correlation-decay theorem applies to the corresponding compact family.
No bound uniform as $p\to\infty$ is asserted or needed.
\end{proposition}
\begin{proof}
The external input is \cite[Theorem 3]{SYZ}, in the precise configuration
class $\widetilde{\mathcal K}$ defined immediately before its Theorem 2
(arXiv:1210.0011v4, pages 6--7). For a fixed finite collection of
$C^{3+\mathrm{Lip}}$ strictly dispersing obstacles, with bounded boundary
derivatives, curvatures bounded away from zero and infinity, positive
minimal flight length, and a common $(t_0,\varphi_0)$ horizon, it gives
\begin{equation}\label{eq:v37-cover-correlation}
 \left|\int f\circ T^m\,g\,d\nu-
                  \int f\,d\nu\int g\,d\nu\right|
        \le C_{f,g}\theta^m,\qquad \theta<1,
\end{equation}
for H\"older $f$ and $1/6$-H\"older $g$. This is a theorem for the
collision map, not just its suspension flow. Its proof is given in
Section 8.2 of that paper. The constant-configuration specialization
of its memory-loss theorem is an equivalent route to
\eqref{eq:v37-cover-correlation}.

We check the geometric hypotheses on each cover. Its planar lift is
exactly the original obstacle array. There are finitely many embedded,
pairwise disjoint obstacles on the cover. The circle and ellipse
boundaries are analytic; on the compact ellipse parameter set their
arclength derivatives through order four and their positive curvatures
have common bounds. The gap is at least $3/50$. Every obstacle contains
the radius-$9/20$ disk. A line avoiding the interiors of those disks
would be a corridor in the smallest circular table, which has none.
Compactness of initial position and direction then gives a bounded
length containing a strictly transverse hit, uniformly in the starting
state. More explicitly, choose a transverse hit for each state, keep
that hit in a neighborhood of the state, and take a finite cover.
The largest hit time and the smallest positive incidence angle in this
finite cover give $t_0$ and $\varphi_0$. Alternatively, first obtaining
a common positive depth of entry into a disk gives the same bounds.
Containment and the common curvature bounds give such a transverse
hit for every member of the compact family. Thus this is the
$(t_0,\varphi_0)$ horizon required in \cite{SYZ}, not merely an upper
bound for almost-everywhere free flights.

The theorem and its local Euclidean collision proof apply on a fixed
flat torus. One can represent the triangular cover by the rectangular
Euclidean cover with periods $(p,0)$ and $(0,p\sqrt3)$ and $2p^2$
obstacles, then pass to its deck quotient. The metric is the inherited
Euclidean metric; no nonisometric affine conjugacy of specular
reflection is performed. Changing the fixed periods only changes the
finite configuration space and the constants in the compactness and
coupling argument. The normalized invariant measure is still
proportional to $\cos\varphi\,dr\,d\varphi$.

Equation~\eqref{eq:v37-cover-correlation} proves mixing first for smooth
functions. To extend it to $f,g\in L^2$, approximate each in $L^2$ by
smooth functions. Invariance and Cauchy--Schwarz bound the replacement
error uniformly in $m$ by
\[
 \|f-f_j\|_2\|g\|_2+\|f_j\|_2\|g-g_j\|_2.
\]
Take $m\to\infty$ before $j\to\infty$. The resulting mixing excludes
all unit-modulus eigenvalues except $1$, and also implies ergodicity.
This verifies exactly the two alternatives in the finite-character
proof: mixing is used when the constant-step character is nonzero;
ergodicity suffices when it is zero. Connectedness of the free region
is only one table hypothesis, not a proof of either conclusion.
\end{proof}

\paragraph{Separated inputs and their uses.}
The hypotheses of Theorem~\ref{thm:v35-action-family-LLT} are now stated
separately in Section~\ref{sec:compact-family-local-principle}.
Local common spaces and unweighted geometric estimates yield the
collision splitting. The fixed preceding-flight partition yields the
analytic multiplier. The action identity and regular matching yield
the three weighted inequalities. Finite-cover mixing excludes physical
peripheral phases. Strong-to-weak parameter continuity gives spectral
uniformity on a fixed compact frequency set. Observation-cell
regularity supplies the varying endpoint tests. Continuity and positive
bounds for the mean roof give the stationary normalization. The
arithmetic step fixes a cover before applying its mixing theorem;
compactness on the base parameter family is used only after pointwise
peripheral exclusion. No summation of estimates over the covers occurs.

\paragraph{Pinned norm references.}
The four unweighted sums in
\eqref{eq:v35-geometric-sums-a}--\eqref{eq:v35-geometric-sums-b} are
\cite[Lemma 3.2(a)--(d), p.\ 12]{DZ}, respectively with the assignment
made just after those displays. The matched graph and stable-Jacobian
estimates are \cite[Lemma 3.3, pp.\ 13--14]{DZ}; its three hypotheses
are regular intermediate stable curves in a common homogeneity strip,
a regular unstable-cone matching foliation, and the intermediate
connector-length bound. The weak, stable and unstable calculations are
\cite[Proposition 4.1, equations (4.1)--(4.3), (4.7)--(4.11),
(4.14), (4.18)]{DZ}. These page and equation numbers refer to
arXiv:1210.1261v1. The piecewise multiplier used here is
\cite[Lemma 3.3(b), p.\ 13]{DPZ}, with its partition conditions (1)--(2)
and the strict inequality $\zeta>\max\{p,\gamma/(1-\gamma)\}$.
The latter numbering refers to arXiv:1902.06850v1, not to a differently
paginated journal copy. These are imported continuum estimates;
source hashes and finite diagnostics do not prove them.
